\chapter{Conclusion}
\label{chapter-5}

\section{Summary}

\par The industrial training program at Active Green Energy Ltd. was an invaluable opportunity to bridge the gap between academic education and practical industry experience. Over the course of the internship, I gained comprehensive exposure to a real business analysis and digital-solutions role, working on two parallel industrial digital product concepts: an Industrial Energy \& Machinery Monitoring System and an IIoT-Based Production Monitoring System. The training enabled me to apply theoretical knowledge in requirements engineering and system design, learn industry-standard tools such as Figma and Jira, and improve my problem-solving and stakeholder-communication skills.

\par This experience enhanced both my technical competencies and professional development. I gained a deeper understanding of how requirements move from a client conversation to a documented specification, then to a design, and finally to a working prototype, as well as how to collaborate with designers, developers, and testers along the way. Overall, the training served as a significant milestone in preparing me for future career opportunities in business analysis and digital solution design.


\section{Challenges During Industrial Training}
\par During the internship, I encountered several challenges that contributed to my learning and growth. Adapting to the company's documentation standards, translating ambiguous client requirements into structured specifications, and coordinating between business and technical stakeholders were all part of the experience.

\par Some specific challenges faced included:
\begin{enumerate}
    \item \textbf{Understanding Incomplete Requirements:} Early client discussions did not always translate directly into clear, testable requirements, requiring follow-up questions and iterative refinement of the SRS/PRD documents.
    \item \textbf{Learning New Tools Quickly:} Becoming proficient with Figma, Jira, and AI-assisted development tools such as Claude Code within a short time frame, while still producing usable output for live projects.
    \item \textbf{Managing a Real Project Delay:} On the IIoT-Based Production Monitoring System project, a delayed factory visit by the IoT engineering team pushed back device setup and the overall timeline, requiring me to help communicate the delay to the client professionally and manage revised expectations.
    \item \textbf{Balancing Two Concurrent Projects:} Dividing attention and documentation effort between two active projects while maintaining consistent quality in both.
    \item \textbf{Validating AI-Assisted Output:} Learning to treat AI-generated prototypes as a starting point requiring careful review against requirements, rather than a finished deliverable.
\end{enumerate}

\section{Scope of the Study}
\par The internship focused on applying and enhancing both academic knowledge and professional skills within the specific domain of business analysis and digital solution design. The following key areas were explored:

\begin{itemize}
    \item \textbf{Understanding Industry Workflows:} Exposure to how requirements, design, and prototyping work in a real industrial technology company.
    \item \textbf{Technical Skill Development:} Practical experience with requirements documentation, system diagramming, UI/UX design, and AI-assisted prototyping tools.
    \item \textbf{Team Collaboration and Communication:} Working with a supervisor, designers, developers, and testers in a coordinated, goal-oriented manner.
    \item \textbf{Professionalism and Ethics:} Learning to communicate difficult news (such as project delays) transparently and professionally.
    \item \textbf{Problem-Solving and Adaptability:} Facing ambiguous requirements and schedule risks, and adapting plans accordingly.
\end{itemize}

\section{Problems and Executions}
\par Throughout the internship, I engaged in tasks that allowed me to apply my academic training in a practical setting -- from requirements analysis coursework to real client requirement-gathering, and from diagramming exercises to production-quality DFD, Use Case, and ER diagrams for live projects.

\par With guidance from my supervisor, I learned how to:
\begin{itemize}
    \item Analyze client requirements and translate them into structured SRS/PRD documentation.
    \item Design system workflows and UI/UX concepts that reflected those requirements accurately.
    \item Use AI-assisted tools to accelerate prototyping while independently validating the results.
    \item Communicate project risks and delays to clients honestly and constructively.
\end{itemize}

\par These experiences helped strengthen my technical foundations in requirements engineering and system design, and taught me how to remain adaptable and solution-oriented when real projects did not go exactly as planned.

\section{Overall Achievement}
\par In conclusion, the internship provided a well-rounded learning experience, helping me grow both as a student and as a future business analyst. I improved my requirements-engineering and UI/UX design skills, gained hands-on experience with AI-assisted development workflows, and learned how to communicate effectively with both technical and non-technical stakeholders. I am particularly proud of successfully carrying both assigned projects from initial requirement gathering through to working prototypes, and of managing the client communication around the IIoT project's delay without it derailing the relationship.

\par If I were to approach the internship again, I would try to establish a lightweight personal log of client feedback and decisions from week one, rather than relying on memory and scattered notes -- a habit I plan to carry into future roles. Overall, the training reinforced my academic learning while introducing me to the practical realities of business analysis, and it has laid a solid foundation for my future career in digital solution design.
